\section{The analytic cumulant reduction}\label{sqt:sec:analytic}
\subsection{Keep the exact remainder on a bounded domain}
Take the analytic logarithm equal to zero at the origin and write
\begin{equation}\label{sqt:eq:h}
h(z)=-\ii z-\log(2-e^{\ii z})+z^2=\sum_{d\geq3}h_dz^d,
\qquad R(x)=\sum_{i,j}h(x_i+x_{n+j}).
\end{equation}
The closest zero of $2-e^{\ii z}$ is $-\ii\log2$, so $h$ is analytic for $|z|<\log2$. Cauchy estimates on a smaller fixed disk give constants independent of $n$ such that, on $|z|\leq C\rho$,
\begin{equation}\label{sqt:eq:derivativebounds}
h'(z)=O(\rho^2),\quad h''(z)=O(\rho),\quad
|h^{(k)}(z)|\leq k!C_0^k\quad(k\geq3).
\end{equation}
The first coefficients are
\begin{equation}\label{sqt:eq:hsix}
h_3=-\ii,\quad h_4=\frac{13}{12},\quad h_5=\frac{5\ii}{4},\quad
h_6=-\frac{541}{360},\quad h_7=-\frac{223\ii}{120},\quad
h_8=\frac{47293}{20160}.
\end{equation}
More generally, if $\genfrac{\{}{\}}{0pt}{}{a}{b}$ is a Stirling number of the second kind,
\begin{equation}\label{sqt:eq:stirling}
h_d=\frac{2\ii^d}{d!}\sum_{j=1}^{d-1}\genfrac{\{}{\}}{0pt}{}{d-1}{j}j!,\qquad d\geq3.
\end{equation}
Indeed the $d$th geometric cumulant is $\sum_{k\geq1}k^{d-1}2^{-k}$, and the falling-factorial expansion of $k^{d-1}$ evaluates that sum as twice the finite Stirling sum. In particular all $h_d/\ii^d$ are rational.

Let $X$ have $2n$ independent $N(0,1/(2n))$ coordinates, and let $\widehat X$ be $X$ conditioned on $\mathcal C_n$. The conditioning event is a product of intervals, so the coordinates remain independent. With $f=R\circ T$,
\begin{equation}\label{sqt:eq:boundedexpectation}
\int_{\mathcal B_n}\mathcal F(x)e^{-(v\cdot x)^2/2}\,\dd x
 =Z_n\PP(X\in\mathcal C_n)\E e^{f(\widehat X)},
\end{equation}
where
\begin{equation}\label{sqt:eq:tailprob}
\PP(X\notin\mathcal C_n)\leq O(n)e^{-n^{2\eps}}
 =O(e^{-c n^{2\eps}}).
\end{equation}
All exponential expectations in this proof use this bounded product domain. Full Gaussians will appear only in fixed polynomial moments and cumulants.

\subsection{Control of every mixed-difference order}
For a cell $e=(i,j)$ write its whitened pair sum as $a_e\cdot y=(Ty)_i+(Ty)_{n+j}$. Its coefficient vector has two entries $1+\gamma/n$ and $2n-2$ entries $\gamma/n$. With $\beta=-\gamma=1-1/\sqrt2$,
\begin{equation}\label{sqt:eq:rowsumbounds}
\sum_{q=1}^{2n}|a_{e,q}|\leq2+2\beta,\qquad
\sum_e|a_{e,q}|\leq(1+\beta)n.
\end{equation}
For a coordinate set $V$ of size $k$, repeated integration of the mixed derivative along the coordinate intervals gives
\[
\Delta_V(f)\leq(2\rho)^k\sum_e
 \sup_{y\in\mathcal C_n}|h^{(k)}(a_e\cdot y)|
 \prod_{q\in V}|a_{e,q}|.
\]
For fixed $q$, an elementary symmetric-sum bound yields
\[
\sum_{\substack{|V|=k\\q\in V}}\prod_{j\in V}|a_{e,j}|
 \leq\frac{|a_{e,q}|}{(k-1)!}\left(\sum_j|a_{e,j}|\right)^{k-1}.
\]
Combining this inequality, \eqref{sqt:eq:derivativebounds} and \eqref{sqt:eq:rowsumbounds}, uniformly for $k\leq2n$, proves
\begin{equation}\label{sqt:eq:Dk}
D_1(f),D_2(f)=O(n\rho^3),\qquad
D_k(f)\leq nk(C\rho)^k\quad(k\geq3).
\end{equation}
These bounds include orders above the selected cumulant truncation.

Fix an integer $p\geq1$, and choose
\begin{equation}\label{sqt:eq:cutofforders}
m=4(p+1),\qquad L=11(p+1),\qquad\alpha=n^{-1/2+4\eps}.
\end{equation}
For each fixed $t\leq m$, the left side of \eqref{sqt:eq:isaevhyp} is $O_m(n\rho^3)$. For example, for $t\geq3$ put $z=C\rho e^{3\alpha/2}$ and use
\[
\sum_{k\geq t}k\binom{k-1}{t-1}z^k
 =\frac{tz^t}{(1-z)^{t+1}}.
\]
The first two terms for $t=1,2$ are controlled by \eqref{sqt:eq:Dk}, and the remaining series is again $O_m(n\rho^3)$. As
$\alpha/(n\rho^3)=n^\eps$, all the hypotheses of Proposition~\ref{sqt:prop:complex} hold for sufficiently large $n$, including $\alpha\leq1/100$. The function is bounded, measurable and complex-valued on the required product domain.

The dimension-adjusted error is
\[
2n\alpha^{m+1}=O\left(n^{1-(1/2-4\eps)(4p+5)}\right)=O(n^{-p}).
\]
At $\eps=1/100$ the exponent is $-1.84p-1.30$, leaving ample slack. Thus
\begin{equation}\label{sqt:eq:finitecumulant}
\E e^{f(\widehat X)}
 =(1+O(n^{-p}))\exp\left(\sum_{r=1}^m\frac{\kappa_r(f(\widehat X))}{r!}\right).
\end{equation}
This is where the order-one cubic phase is controlled. Expanding that phase as if it were a uniformly small perturbation would not justify \eqref{sqt:eq:finitecumulant}.

\subsection{Replace only the finite cumulants by polynomials}
Let
\[
R_L(x)=\sum_{d=3}^L h_d\sum_{i,j}(x_i+x_{n+j})^d.
\]
On $\mathcal B_n$,
\begin{equation}\label{sqt:eq:RLbounds}
|R|+|R_L|=O(n^2\rho^3),\qquad |R-R_L|=O(n^2\rho^{L+1}).
\end{equation}
The moment-partition formula for a cumulant, with one factor subtracted at a time, bounds the difference of each order-$r$ cumulant by
\[
C_r\|R-R_L\|_\infty
 (1+\|R\|_\infty+\|R_L\|_\infty)^{r-1}.
\]
For $r\leq m$ the exponent of $n$ in this bound is at most
\begin{equation}\label{sqt:eq:replacementexponent}
1+\frac m2-\frac L2+\eps(L+3m-2)
 \leq-\frac{327}{100}p-\frac{229}{100}<-p.
\end{equation}
Hence replacing $R$ by $R_L$ inside the finite cumulant sum costs $O(n^{-p})$.

For fixed $r,L$, the polynomial $R_L(TX)^r$ has fixed degree; the number of its terms and the magnitude of its coefficients are bounded by powers of $n$. Fixed Gaussian moments are likewise polynomially bounded. By Cauchy--Schwarz and \eqref{sqt:eq:tailprob}, conditioning on $\mathcal C_n$ changes any required moment by at most a polynomial in $n$ times $e^{-c n^{2\eps}}$. Applying the moment-partition formula again gives
\begin{equation}\label{sqt:eq:fullgaussreplace}
\kappa_r(R_L(T\widehat X))=\kappa_r(R_L(TX))+O(n^{-p}),\qquad r\leq m.
\end{equation}
No unbounded exponential of $R_L$ has been introduced in this step.

Set $Y=TX$ and $H_{p,n}=\sum_{r=1}^m\kappa_r(R_L(Y))/r!$. The diagram bound in the next section proves that $H_{p,n}$ is real and $O(1)$. Therefore $e^{H_{p,n}}$ is bounded above and away from zero. Combining \eqref{sqt:eq:transfer}, \eqref{sqt:eq:prefactor}, \eqref{sqt:eq:boundedexpectation}, and \eqref{sqt:eq:finitecumulant}--\eqref{sqt:eq:fullgaussreplace} gives
\begin{equation}\label{sqt:eq:analyticreduction}
a_n=G_n\exp(H_{p,n})(1+O(n^{-p})).
\end{equation}
The exponentially small absolute localization error is now a relative negligible error. Positivity and the real logarithm in subsequent formulas are justified by the real bounded exponent and a factor tending to $1$; no nonvanishing assumption has been added to Isaev's theorem.
